%% notes-tex/database/ecoli-datasets/sections/datasets.tex
%% [[experiments.database.ecoli-datasets]]
%% https://github.com/Mjvolk3/torchcell/tree/main/notes-tex/database/ecoli-datasets/sections/datasets.tex
%% SOURCE: hand-authored prose; every number quoted here is a macro from
%% tables/numbers.tex, generated by
%% experiments/database/scripts/build_ecoli_genslm_benchmark_table.py off the snapshot
%% written by experiments/database/scripts/build_bacteria_datasets_table.py

\section{The \emph{E. coli} datasets\secstatus{todo}}\label{sec:datasets}

Table~\ref{tab:ecoli-datasets} lists the \ecoliDatasetCount{} datasets in
\ecoliSectionCount{} phenotype families. \term{Genotypes} and \term{Env} are the
distinct stored genotype and environment payloads, \term{Instances} the record count,
and \term{Signal} the gzip size in bits of the concatenated stored instances, the same
relative proxy the yeast table uses. A dataset whose strains are all the one wild type,
such as a proteome panel across growth conditions, shows one genotype and as many
environments as conditions; a deletion screen shows one genotype per strain.

Two columns are specific to this document. The \term{sequence-only fit} tier says how
well a model that sees only sequence embeddings fits the dataset, and the \term{label
unit} names what one label attaches to once the cell graph is dropped. Both are derived
by the probe script from where the label sits in the cell graph (a whole-strain readout,
a node, a gene pair, a metabolite) and from the modal number of edits in a sample of
400 genotypes per store. The tiers are structural: \emph{best} when the label sits on
one gene, one gene pair or triple, or one inserted construct, so that a per-gene or
per-pair representation is the whole input; \emph{good} when the label needs the readout
entity crossed with the perturbed genes or with the condition, which a pair of vectors
still carries; \emph{weak} for a multi-edit strain where only its heterologous genes have
a sequence to stand on; \emph{no} for a strain readout with no single edit to attach a
sequence to. Table~\ref{tab:tiers} sums the datasets and instances per tier. A tier is a
reading of the genotype, not a model result; it says what a sequence-only model could
be given, not what it would score.

%% GENERATED FILE -- do not hand-edit.
%% SOURCE: experiments/database/scripts/build_ecoli_genslm_benchmark_table.py
%% off experiments/database/results/pre-build/2026-10-21/ecoli_datasets.json
\begin{table}[H]
\centering
\footnotesize
\setlength{\tabcolsep}{4pt}
\caption{\emph{E. coli} datasets and instances per sequence-only fit tier, as \texttt{sequence\_fit} in the probe script derives it from where the label sits and the modal edit count of the genotype sample.}
\label{tab:tiers}
\begin{tabular}{@{}l r r >{\raggedright\arraybackslash}p{88mm}@{}}
\toprule
Sequence-only fit & Datasets & Instances & What the tier means \\
\midrule
best & 35 & 639,394 & the label sits on one gene, a gene pair or triple, or one inserted construct \\
good & 29 & 3,921,964 & readout entity x perturbed gene(s), or gene x condition \\
weak & 1 & 72 & a multi-edit strain, only its heterologous genes carry a sequence \\
no & 3 & 435 & a natural genome, or a whole-strain readout with no single edit \\
\midrule
\textbf{Total} & \textbf{68} & \textbf{4,561,865} &  \\
\bottomrule
\end{tabular}
\end{table}


\clearpage
%% GENERATED FILE -- do not hand-edit.
%% SOURCE: experiments/database/scripts/build_ecoli_genslm_benchmark_table.py
%% off experiments/database/results/pre-build/2026-10-21/ecoli_datasets.json
\begin{landscape}
\begingroup
\scriptsize
\setlength{\tabcolsep}{3pt}
\begin{longtable}{@{}l r r r >{\raggedright\arraybackslash}p{34mm} l l r l >{\raggedright\arraybackslash}p{40mm}@{}}
\caption{The \emph{E. coli} datasets served from the dev tree, snapshot 2026-10-21, grouped by phenotype family. \emph{Genotypes} and \emph{Env} count distinct stored genotype and environment payloads; \emph{Instances} is the record count; \emph{Signal} the gzip size in bits of the concatenated stored instances (perturbation, environment and phenotype). \emph{Seq.\ fit} is the sequence-only fit tier of Table~\ref{tab:tiers} and \emph{Label unit} what one label attaches to without the cell graph, both read from the genotype of 400 evenly spaced records per store.}\label{tab:ecoli-datasets}\\
\toprule
Dataset & Genotypes & Env & Instances & Phenotype & Shape & Role & Signal (bits) & Seq. fit & Label unit \\
\midrule
\endfirsthead
\multicolumn{10}{@{}l}{\footnotesize\emph{Table~\ref{tab:ecoli-datasets}, continued}}\\
\toprule
Dataset & Genotypes & Env & Instances & Phenotype & Shape & Role & Signal (bits) & Seq. fit & Label unit \\
\midrule
\endhead
\multicolumn{10}{@{}l}{\textbf{Environmental / chemogenomic response}} \\
CrisprPineneToleranceNiu2019 & 54 & 1 & 54 & environment\_\allowbreak{}response (log2\_\allowbreak{}ratio) & scalar & global & $4.1\times10^{4}$ & best & 1 gene (1-6 edits) \\
CrispriChemgenChoe2025 & 39,291 & 12 & 470,502 & environment\_\allowbreak{}response (log2\_\allowbreak{}ratio) & scalar & global & $2.4\times10^{8}$ & good & 1 gene x condition \\
CrispriGeneFitnessWang2018 & 13,641 & 5 & 20,158 & environment\_\allowbreak{}response (log2\_\allowbreak{}ratio) & scalar & global & $1.5\times10^{7}$ & good & 1 gene x condition (1-7 edits) \\
CrispriGuideFfaEnrichmentFang2025 & 15,711 & 1 & 15,711 & environment\_\allowbreak{}response (log2\_\allowbreak{}ratio) & scalar & global & $7.3\times10^{6}$ & best & 1 gene (1-5 edits) \\
CrispriGuideFitnessWang2018 & 53,203 & 5 & 240,481 & environment\_\allowbreak{}response (log2\_\allowbreak{}ratio) & scalar & global & $1.6\times10^{8}$ & good & 1 gene x condition \\
CrispriKnockdownCui2018 & 71,696 & 1 & 143,392 & environment\_\allowbreak{}response (log2\_\allowbreak{}ratio) & scalar & global & $7.3\times10^{7}$ & best & 1 gene (1-7 edits) \\
CrispriScreenRousset2018 & 40,592 & 4 & 92,252 & environment\_\allowbreak{}response (log2\_\allowbreak{}ratio) & scalar & global & $4.2\times10^{7}$ & good & 1 gene x condition \\
DoublingTimeCaglar2017 & 1 & 16 & 55 & environment\_\allowbreak{}response (growth\_\allowbreak{}rate) & scalar & global & $3.2\times10^{4}$ & no & unperturbed strain readout x condition \\
EnvChemgenGirgis2009 & 3,976 & 17 & 66,287 & environment\_\allowbreak{}response (z\_\allowbreak{}score) & scalar & global & $3.7\times10^{7}$ & good & 1 gene x condition \\
EnvChemgenNichols2011 & 3,967 & 324 & 1,286,962 & environment\_\allowbreak{}response (z\_\allowbreak{}score) & scalar & global & $4.1\times10^{8}$ & good & 1 gene x condition \\
EnvChemgenShiver2016 & 3,846 & 51 & 210,998 & environment\_\allowbreak{}response (z\_\allowbreak{}score) & scalar & global & $9.9\times10^{7}$ & good & 1 gene x condition \\
EnvChemgenWang2015 & 46 & 1 & 46 & environment\_\allowbreak{}response (log2\_\allowbreak{}ratio) & scalar & global & $3.0\times10^{4}$ & best & 1 gene \\
EnvChemgenWang2024 & 4,419 & 6 & 26,514 & environment\_\allowbreak{}response (log2\_\allowbreak{}ratio) & scalar & global & $9.3\times10^{6}$ & good & 1 gene x condition \\
FitnessBrowserEcoli & 3,789 & 3 & 22,734 & environment\_\allowbreak{}response (log2\_\allowbreak{}ratio) & scalar & global & $1.1\times10^{7}$ & good & 1 gene x condition \\
GrowthRateLamoureux2023 & 21 & 54 & 354 & environment\_\allowbreak{}response (growth\_\allowbreak{}rate) & scalar & global & $2.1\times10^{5}$ & no & unperturbed strain readout x condition (0-1 edits) \\
GrowthRateS23Schmidt2016 & 1 & 22 & 26 & environment\_\allowbreak{}response (growth\_\allowbreak{}rate) & scalar & global & $2.0\times10^{4}$ & no & unperturbed strain readout x condition \\
GrowthYieldBaba2006 & 3,912 & 3 & 11,735 & environment\_\allowbreak{}response (optical\_\allowbreak{}density) & scalar & global & $3.3\times10^{6}$ & good & 1 gene x condition \\
MismatchCrispriFitnessHawkins2020 & 24,186 & 1 & 24,186 & environment\_\allowbreak{}response (relative\_\allowbreak{}growth\_\allowbreak{}rate) & scalar & global & $1.4\times10^{7}$ & best & 1 gene \\
PhageRbTnseqMutalik2020 & 3,716 & 71 & 287,815 & environment\_\allowbreak{}response (log2\_\allowbreak{}ratio) & scalar & global & $2.1\times10^{8}$ & good & 1 gene x condition \\
RbTnseqPrice2018Ecoli & 3,789 & 104 & 598,662 & environment\_\allowbreak{}response (log2\_\allowbreak{}ratio) & scalar & global & $4.0\times10^{8}$ & good & 1 gene x condition \\
RbTnseqTeteneva2024 & 3,691 & 18 & 66,162 & environment\_\allowbreak{}response (log2\_\allowbreak{}ratio) & scalar & global & $2.7\times10^{7}$ & good & 1 gene x condition \\
RbTnseqWetmore2015 & 3,646 & 85 & 353,662 & environment\_\allowbreak{}response (log2\_\allowbreak{}ratio) & scalar & global & $2.4\times10^{8}$ & good & 1 gene x condition \\
\addlinespace
\multicolumn{10}{@{}l}{\textbf{Expression (RNA-seq)}} \\
MrnaFractionBalakrishnan2022 & 3 & 14 & 28 & mrna\_\allowbreak{}number\_\allowbreak{}fraction (rnaseq\_\allowbreak{}mrna\_\allowbreak{}number\_\allowbreak{}fraction) & vector (4342) & node & $7.2\times10^{6}$ & best & 1 gene, unperturbed strain x condition (0-2 edits) \\
RnaseqCaglar2017 & 1 & 57 & 152 & expression\_\allowbreak{}tpm (rnaseq\_\allowbreak{}tpm) & vector (4196) & node & $6.7\times10^{7}$ & best & 1 gene, unperturbed strain x condition \\
RnaseqGuzman2018 & 7 & 2 & 14 & expression\_\allowbreak{}tpm (rnaseq\_\allowbreak{}tpm) & vector (4329) & node & $4.5\times10^{6}$ & good & readout gene x 3 perturbed gene(s) x condition (0-5 edits) \\
RnaseqLamoureux2023 & 74 & 119 & 1,035 & expression\_\allowbreak{}tpm (rnaseq\_\allowbreak{}tpm) & vector (4257) & node & $5.8\times10^{8}$ & best & 1 gene, unperturbed strain x condition (0-3 edits) \\
RnaseqPublicK12Lamoureux2023 & 72 & 198 & 1,675 & expression\_\allowbreak{}tpm (rnaseq\_\allowbreak{}tpm\_\allowbreak{}from\_\allowbreak{}released\_\allowbreak{}counts) & vector (4355) & node & $9.5\times10^{8}$ & best & 1 gene, unperturbed strain x condition (0-4 edits) \\
\addlinespace
\multicolumn{10}{@{}l}{\textbf{Fitness}} \\
CarbonSourceTong2020 & 3,714 & 30 & 111,420 & fitness & scalar & global & $4.8\times10^{7}$ & good & 1 gene x condition \\
CrispriCrossRachwalski2024 & 14,209 & 15 & 48,900 & fitness & scalar & global & $1.9\times10^{7}$ & good & gene pair, strain readout x condition (1-2 edits) \\
GrowthAucRapp2026 & 1,515 & 1 & 1,515 & fitness & scalar & global & $7.6\times10^{5}$ & best & 1 gene \\
GrowthRateCampos2018 & 3,664 & 1 & 3,664 & fitness & scalar & global & $1.1\times10^{6}$ & best & 1 gene \\
GrowthRateChoe2019 & 6 & 1 & 8 & fitness & scalar & global & $1.1\times10^{4}$ & best & 1 gene (0-21 edits) \\
GrowthRateSchmidt2016 & 3 & 2 & 6 & fitness & scalar & global & $8.5\times10^{3}$ & good & 1 gene x condition \\
GrowthWang2015 & 46 & 1 & 46 & fitness & scalar & global & $2.2\times10^{4}$ & best & 1 gene \\
TranscriptionFactorKnockoutChoe2019 & 2 & 1 & 2 & fitness & scalar & global & $5.8\times10^{3}$ & best & 1 gene \\
\addlinespace
\multicolumn{10}{@{}l}{\textbf{Flux}} \\
FluxIshii2007 & 25 & 5 & 28 & net\_\allowbreak{}flux (c13\_\allowbreak{}mfa\_\allowbreak{}net\_\allowbreak{}flux\_\allowbreak{}percent\_\allowbreak{}of\_\allowbreak{}specific\_\allowbreak{}glucose\_\allowbreak{}uptake) & vector (42) & bipartite node & $6.0\times10^{4}$ & good & reaction x 1 perturbed gene(s) x condition (0-1 edits) \\
FluxLi2021 & 3 & 1 & 3 & net\_\allowbreak{}flux (c13\_\allowbreak{}mfa\_\allowbreak{}net\_\allowbreak{}flux\_\allowbreak{}mmol\_\allowbreak{}per\_\allowbreak{}gdw\_\allowbreak{}per\_\allowbreak{}h) & vector (99) & bipartite node & $5.7\times10^{4}$ & good & reaction x 1 perturbed gene(s) \\
\addlinespace
\multicolumn{10}{@{}l}{\textbf{Gene interaction}} \\
GeneInteractionBabu2014 & 42,703 & 1 & 42,705 & gene\_\allowbreak{}interaction & scalar & edge & $1.2\times10^{7}$ & best & gene pair \\
GeneInteractionButland2008 & 314,769 & 1 & 314,769 & gene\_\allowbreak{}interaction & scalar & edge & $8.2\times10^{7}$ & best & gene pair \\
\addlinespace
\multicolumn{10}{@{}l}{\textbf{Metabolite}} \\
ExometaboliteBrunk2016 & 9 & 28 & 126 & metabolite\_\allowbreak{}level (hplc\_\allowbreak{}extracellular\_\allowbreak{}concentration\_\allowbreak{}g\_\allowbreak{}per\_\allowbreak{}l) & vector (6) & bipartite node & $2.3\times10^{5}$ & good & metabolite x 9 perturbed gene(s) x condition (0-10 edits) \\
MetaboliteIntensityRapp2026 & 407 & 1 & 407 & metabolite\_\allowbreak{}level (fi\_\allowbreak{}ms\_\allowbreak{}annotated\_\allowbreak{}feature\_\allowbreak{}absolute\_\allowbreak{}intensity\_\allowbreak{}mean\_\allowbreak{}of\_\allowbreak{}2\_\allowbreak{}plates) & vector (2) & bipartite node & $3.9\times10^{5}$ & good & metabolite x 1 perturbed gene(s) \\
MetabolomeBrunk2016 & 9 & 26 & 117 & metabolite\_\allowbreak{}level (lc\_\allowbreak{}ms\_\allowbreak{}concentration\_\allowbreak{}um) & vector (51) & bipartite node & $4.5\times10^{5}$ & good & metabolite x 9 perturbed gene(s) x condition (0-10 edits) \\
MetabolomeFuhrer2017 & 3,806 & 1 & 3,806 & metabolite\_\allowbreak{}level (fia\_\allowbreak{}tof\_\allowbreak{}ms\_\allowbreak{}ion\_\allowbreak{}modified\_\allowbreak{}z\_\allowbreak{}score) & vector (7534) & bipartite node & $1.9\times10^{9}$ & good & metabolite x 1 perturbed gene(s) \\
MetabolomeIshii2007 & 25 & 5 & 28 & metabolite\_\allowbreak{}level (ce\_\allowbreak{}tofms\_\allowbreak{}intracellular\_\allowbreak{}concentration\_\allowbreak{}mm) & vector (130) & bipartite node & $4.9\times10^{5}$ & good & metabolite x 1 perturbed gene(s) x condition (0-1 edits) \\
MetabolomeRapp2026 & 1,497 & 1 & 1,497 & metabolite\_\allowbreak{}level (fi\_\allowbreak{}ms\_\allowbreak{}iml1515\_\allowbreak{}feature\_\allowbreak{}fold\_\allowbreak{}change\_\allowbreak{}vs\_\allowbreak{}batch\_\allowbreak{}median\_\allowbreak{}mean\_\allowbreak{}of\_\allowbreak{}2\_\allowbreak{}plates) & vector (1321) & bipartite node & $5.6\times10^{8}$ & good & metabolite x 1 perturbed gene(s) \\
MetabolomeSchastnaya2021 & 103 & 5 & 200 & metabolite\_\allowbreak{}level (fia\_\allowbreak{}tof\_\allowbreak{}ms\_\allowbreak{}ion\_\allowbreak{}log2\_\allowbreak{}fold\_\allowbreak{}change\_\allowbreak{}vs\_\allowbreak{}wild\_\allowbreak{}type) & vector (284) & bipartite node & $3.6\times10^{6}$ & good & metabolite x 1 perturbed gene(s) x condition \\
TargetedMetabolomeRapp2026 & 407 & 1 & 407 & metabolite\_\allowbreak{}level (lc\_\allowbreak{}msms\_\allowbreak{}eic\_\allowbreak{}peak\_\allowbreak{}height\_\allowbreak{}fold\_\allowbreak{}change\_\allowbreak{}vs\_\allowbreak{}library\_\allowbreak{}median) & vector (2) & bipartite node & $3.6\times10^{5}$ & good & metabolite x 1 perturbed gene(s) \\
\addlinespace
\multicolumn{10}{@{}l}{\textbf{Morphology}} \\
MorphologyCampos2018 & 3,664 & 1 & 3,664 & morphology & vector (15) & global & $8.5\times10^{6}$ & best & 1 gene \\
\addlinespace
\multicolumn{10}{@{}l}{\textbf{Phosphoproteome}} \\
PhosphositeIntensityPotel2018 & 1 & 2 & 12 & phosphosite\_\allowbreak{}intensity (fe3\_\allowbreak{}imac\_\allowbreak{}lc\_\allowbreak{}msms\_\allowbreak{}maxquant\_\allowbreak{}summed\_\allowbreak{}xic) & vector (989) & node & $1.3\times10^{7}$ & best & 1 gene, unperturbed strain x condition \\
\addlinespace
\multicolumn{10}{@{}l}{\textbf{Product titer}} \\
BiofuelTiterBrunk2016 & 8 & 18 & 72 & titer & scalar & global & $1.5\times10^{5}$ & weak & set of 9 edits incl. heterologous genes (7-10 edits) \\
IsopentenolTiterFoo2014 & 9 & 1 & 9 & titer & scalar & global & $1.7\times10^{4}$ & best & 1 heterologous construct \\
PineneFoldChangeNiu2019 & 40 & 1 & 40 & titer & scalar & global & $2.3\times10^{4}$ & best & 1 gene (1-6 edits) \\
\addlinespace
\multicolumn{10}{@{}l}{\textbf{Promoter activity}} \\
PromoterReporterMohiuddin2022 & 1,809 & 27 & 69,480 & promoter\_\allowbreak{}activity & scalar & node & $3.3\times10^{7}$ & best & inserted DNA (promoter or reporter) x condition \\
\addlinespace
\multicolumn{10}{@{}l}{\textbf{Protein abundance}} \\
ProteinFoldChangeCaglar2017 & 1 & 8 & 8 & protein\_\allowbreak{}fold\_\allowbreak{}change (deseq2\_\allowbreak{}wald\_\allowbreak{}log2\_\allowbreak{}fold\_\allowbreak{}change\_\allowbreak{}design\_\allowbreak{}batch\_\allowbreak{}plus\_\allowbreak{}condition) & vector (4180) & node & $1.1\times10^{7}$ & best & 1 gene, unperturbed strain x condition \\
ProteomeBrunk2016 & 9 & 18 & 81 & protein\_\allowbreak{}abundance (srm\_\allowbreak{}protein\_\allowbreak{}peak\_\allowbreak{}area) & vector (68) & node & $4.2\times10^{5}$ & good & readout gene x 9 perturbed gene(s) x condition (0-10 edits) \\
ProteomeCaglar2017 & 1 & 40 & 105 & protein\_\allowbreak{}abundance (lcmsms\_\allowbreak{}spectral\_\allowbreak{}count\_\allowbreak{}deseq2\_\allowbreak{}size\_\allowbreak{}factor\_\allowbreak{}normalized) & vector (4196) & node & $2.5\times10^{7}$ & best & 1 gene, unperturbed strain x condition \\
ProteomeIshii2007 & 25 & 5 & 28 & protein\_\allowbreak{}abundance (lc\_\allowbreak{}ms\_\allowbreak{}ms\_\allowbreak{}absolute\_\allowbreak{}mg\_\allowbreak{}protein\_\allowbreak{}per\_\allowbreak{}g\_\allowbreak{}dry\_\allowbreak{}cell\_\allowbreak{}weight) & vector (57) & node & $2.9\times10^{5}$ & best & 1 gene, its own label x condition (0-1 edits) \\
ProteomeMori2021 & 1 & 1 & 7 & protein\_\allowbreak{}abundance (absolute\_\allowbreak{}protein\_\allowbreak{}mass\_\allowbreak{}fraction\_\allowbreak{}xtop\_\allowbreak{}dia\_\allowbreak{}swath\_\allowbreak{}riboprofiling\_\allowbreak{}scaled) & vector (1897) & node & $1.1\times10^{6}$ & best & 1 gene, unperturbed strain \\
ProteomeSchmidt2016 & 1 & 21 & 21 & protein\_\allowbreak{}abundance (absolute\_\allowbreak{}protein\_\allowbreak{}copies\_\allowbreak{}per\_\allowbreak{}cell\_\allowbreak{}label\_\allowbreak{}free\_\allowbreak{}sid\_\allowbreak{}anchored) & vector (2359) & node & $1.1\times10^{7}$ & best & 1 gene, unperturbed strain x condition \\
ProteomeSrmSet1Schmidt2016 & 1 & 18 & 18 & protein\_\allowbreak{}abundance (absolute\_\allowbreak{}protein\_\allowbreak{}copies\_\allowbreak{}per\_\allowbreak{}cell\_\allowbreak{}srm\_\allowbreak{}sid\_\allowbreak{}dataset1\_\allowbreak{}technical\_\allowbreak{}sd) & vector (41) & node & $1.5\times10^{5}$ & best & 1 gene, unperturbed strain x condition \\
ProteomeSrmSet2Schmidt2016 & 1 & 21 & 21 & protein\_\allowbreak{}abundance (absolute\_\allowbreak{}protein\_\allowbreak{}copies\_\allowbreak{}per\_\allowbreak{}cell\_\allowbreak{}srm\_\allowbreak{}sid\_\allowbreak{}dataset2\_\allowbreak{}biological\_\allowbreak{}sd) & vector (31) & node & $1.3\times10^{5}$ & best & 1 gene, unperturbed strain x condition \\
\addlinespace
\multicolumn{10}{@{}l}{\textbf{Protein synthesis / turnover}} \\
ProteinSynthesisRateLi2014 & 1 & 3 & 3 & synthesis\_\allowbreak{}rate (ribosome\_\allowbreak{}profiling\_\allowbreak{}footprint\_\allowbreak{}density\_\allowbreak{}absolute\_\allowbreak{}synthesis\_\allowbreak{}rate) & vector (3041) & node & $3.7\times10^{5}$ & best & 1 gene, unperturbed strain x condition \\
ProteinTurnoverGupta2024 & 6 & 8 & 13 & degradation\_\allowbreak{}rate (n15\_\allowbreak{}ammonium\_\allowbreak{}tmtproc\_\allowbreak{}total\_\allowbreak{}turnover\_\allowbreak{}rate\_\allowbreak{}per\_\allowbreak{}hour) & vector (2553) & node & $1.1\times10^{7}$ & best & 1 gene, unperturbed strain x condition (0-3 edits) \\
\addlinespace
\multicolumn{10}{@{}l}{\textbf{Viability}} \\
GeneEssentialityBaba2006 & 4,288 & 1 & 4,288 & is\_\allowbreak{}essential & scalar & node & $1.8\times10^{6}$ & best & 1 gene, its own label \\
GeneEssentialityGerdes2003 & 3,739 & 1 & 3,739 & is\_\allowbreak{}essential & scalar & node & $2.0\times10^{6}$ & best & 1 gene, its own label \\
GeneEssentialityGoodall2018 & 8,309 & 2 & 8,309 & is\_\allowbreak{}essential & scalar & node & $5.2\times10^{6}$ & best & 1 gene, its own label x condition \\
GeneEssentialityPrice2018Ecoli & 324 & 1 & 324 & is\_\allowbreak{}essential & scalar & node & $2.6\times10^{5}$ & best & 1 gene, its own label \\
GradedGeneEssentialityGoodall2018 & 317 & 2 & 317 & essentiality\_\allowbreak{}call & scalar & node & $2.2\times10^{5}$ & best & 1 gene, its own label x condition \\
\midrule
\textbf{Total (68 datasets)} &  &  & \textbf{4,561,865} &  &  &  & \textbf{$6.4\times10^{9}$} &  &  \\
\bottomrule
\end{longtable}
\endgroup
\end{landscape}

